% !PRINTED 73
\weberpage{73}

Formula (2), applied to this system of functions, yields the
following, frequently used table:

\weq{8}{
\begin{aligned}
\vartheta_{00}(v) &= \vartheta_{01}\!\left(v+\tfrac{1}{2}\right) = \varepsilon\,\vartheta_{10}\!\left(v+\tfrac{\omega}{2}\right) = \varepsilon\,\vartheta_{11}\!\left(v+\tfrac{1+\omega}{2}\right),\\
\vartheta_{01}(v) &= \vartheta_{00}\!\left(v+\tfrac{1}{2}\right) = -i\varepsilon\,\vartheta_{11}\!\left(v+\tfrac{\omega}{2}\right) = i\varepsilon\,\vartheta_{10}\!\left(v+\tfrac{1+\omega}{2}\right),\\
\vartheta_{10}(v) &= \vartheta_{11}\!\left(v+\tfrac{1}{2}\right) = \varepsilon\,\vartheta_{00}\!\left(v+\tfrac{\omega}{2}\right) = \varepsilon\,\vartheta_{01}\!\left(v+\tfrac{1+\omega}{2}\right),\\
\vartheta_{11}(v) &= -\vartheta_{10}\!\left(v+\tfrac{1}{2}\right) = -i\varepsilon\,\vartheta_{01}\!\left(v+\tfrac{\omega}{2}\right) = -i\varepsilon\,\vartheta_{11}\!\left(v+\tfrac{1+\omega}{2}\right),
\end{aligned}
}

in which, for brevity,

\[
\varepsilon = e^{\frac{\pi i \omega}{4} + \pi i v}
\]

has been set.

According to (3), (4) we have

\weq{9}{
\begin{aligned}
\vartheta_{00}(-v) &= \vartheta_{00}(v), & \vartheta_{01}(-v) &= \vartheta_{01}(v),\\
\vartheta_{10}(-v) &= \vartheta_{10}(v), & \vartheta_{11}(-v) &= -\vartheta_{11}(v),
\end{aligned}
}

that is, \gesperrt{$\vartheta_{00}(v)$, $\vartheta_{01}(v)$, $\vartheta_{10}(v)$ are even functions,
$\vartheta_{11}(v)$ is an odd function.}

According to § 17 (12) the four functions (7) vanish, respectively,
for the following values of the argument

\[
\frac{1+\omega}{2},\quad \frac{\omega}{2},\quad \frac{1}{2},\quad 0,
\]

hence for certain half periods. Of importance are the values which
all the functions (7) assume for these values of the argument, and
these can be reduced, by means of the table (8), to the three constants:

\[
\vartheta_{00}(0) = \vartheta_{00},\quad \vartheta_{01}(0) = \vartheta_{01},\quad \vartheta_{10}(0) = \vartheta_{10}
\]

One thus obtains from (8) the following system of
formulae:

\weq{10}{
\begin{aligned}
\vartheta_{00} &= \vartheta_{01}\!\left(\tfrac{1}{2}\right) = \varepsilon_0\,\vartheta_{10}\!\left(\tfrac{\omega}{2}\right) = \varepsilon_0\,\vartheta_{11}\!\left(\tfrac{1+\omega}{2}\right),\\
\vartheta_{01} &= \vartheta_{00}\!\left(\tfrac{1}{2}\right) = -i\varepsilon_0\,\vartheta_{11}\!\left(\tfrac{\omega}{2}\right) = i\varepsilon_0\,\vartheta_{10}\!\left(\tfrac{1+\omega}{2}\right),\\
\vartheta_{10} &= \vartheta_{11}\!\left(\tfrac{1}{2}\right) = \varepsilon_0\,\vartheta_{00}\!\left(\tfrac{\omega}{2}\right) = \varepsilon_0\,\vartheta_{01}\!\left(\tfrac{1+\omega}{2}\right),
\end{aligned}
}
% [sic] German source repeats theta_11 at the far right of row four in (8).
%CONT>>
